% English translation of checked_pages/esq_119_023.tex; source: cote 119, site 23, PDF page 24, typescript page 21.
% AI-translated by GPT-6 Astra (OpenAI, ChatGPT Pro mode).
\scanpage{119/23}
\mspage{21}

My sporadic reflection on this question over several years
was confined to bringing out a few simple basic principles, at
first directing my attention to \uline{pinned} regular
polyhedra,\tnote{TN-023-01}{``Pinned'' renders ``épinglé'', continuing the source's language of fixing a polyhedron by choosing a frame.}
which reduces the necessary conceptual baggage to a minimum
and practically eliminates any moderately delicate rationality
questions. For such a polyhedron one finds a canonical basis (or
frame) of the ambient affine or projective space, in
such a way that the operations of the cartographic group $\underline{C}_n$, generated
by the fundamental reflections $\sigma_i$ $(0 \leq i \leq n)$, are written in it
by universal formulas in terms of $n$ parameters
$\alpha_1,\ldots,\alpha_n$ , which geometrically are interpreted as twice
the cosines of the ``fundamental angles'' of the polyhedron. The polyhedron
is reconstructed from this action and the affine or
projective flag associated with the chosen basis, by transforming that flag
by all the elements of the group generated by the fundamental
reflections. Thus the ``universal'' pinned n-polyhedron is
canonically defined over the polynomial ring in $n$ indeterminates
\[ \mathbb{Z}[\underline{\alpha}_1,\ldots,\underline{\alpha}_n] , \]
its specializations over arbitrary base fields $k$ (through
values $\alpha_i \in k$ assigned to the indeterminates $\underline{\alpha}_i$) giving
regular polyhedra corresponding to various combinatorial
types. In this game there is no question of restricting oneself to
finite regular polyhedra, or even to regular polyhedra whose
facets have finite order, i.e. for which the parameters
$\alpha_i$ are roots of suitable ``semicyclotomic'' equations
expressing that the ``fundamental angles'' (in the
case where the base field is $\mathbb{R}$) are commensurable with $2\pi$. Already
when $n = 1$, perhaps the most interesting regular polygon
of all (morally, the regular polygon with just one side !)
is the one corresponding to $\alpha = 2$, giving rise to a circumscribed
parabolic conic, i.e. tangent to the line at infinity.

The finite case is the one in which the group generated by the fundamental
reflections, which is also the automorphism group of the

\margin{23}

